\documentclass[11pt, a4paper]{report} % Classe: report (per capitoli) o article

% Pacchetti fondamentali
\usepackage[utf8]{inputenc} % Codifica caratteri
\usepackage[T1]{fontenc}    % Codifica font
\usepackage[italian]{babel} % Lingua
\usepackage{graphicx}       % Inserimento immagini
\usepackage{amsmath}        % Formule matematiche
\usepackage{hyperref} 


% Dati della relazione
\title{Relazione Progetto Business Intelligence Per i Servizi Finanziari 25/26}
\author{Luca Giani [909374]}
\date{\today}
\begin{document}

\maketitle % Crea il frontespizio
\tableofcontents % Indice automatico

\chapter{Introduzione}

\section{Titoli scelti}

I 6 titoli appartenenti all'indice S\&P 500 scelti per l'analisi sono stati:
\begin{enumerate}
    \item Settore Tecnologico: 
    \begin{itemize}
        \item NVIDIA (NVDA): leader nel settore dei semiconduttori e dell'intelligenza artificiale, con una forte crescita negli ultimi anni. È stata scelta visto il suo ruolo chiave nel'AI che è un settore in forte espansione, visto il rally che ha avuto negli ultimi anni e la recente notizia di 1 Trilione di ricavi cumulativi entro la fine del 2027.
        \item Microsoft (MSFT): gigante tecnologico con una vasta gamma di prodotti e servizi, tra cui software, cloud computing e intelligenza artificiale. È stato scelto proprio per la sua capcità di reinventarsi durante gli anni e rimanere per anni all'interno dei "magnificent-seven".
    \end{itemize}
        
    \item Settore Energetico: 
    \begin{itemize}
        \item ExxonMobil (XOM): azienda petrolifera leader mondiale, con una forte presenza in tutto il mondo. È stata scelta per la sua posizione di leadership nel settore energetico, essendo la compagnia americana con la maggiore capitalizzazione di mercato e perchè a differenza delle altre compagnie è più caute riguardo gli investimenti nelle energie rinnovabili.
        \item Chevron (CVX): azienda petrolifera americana, con una forte presenza negli Stati Uniti e in altri paesi. È stata scelta per il suo orientatmento pìù proiettato verso le energie rinnovabili rispetto a ExxonMobil, visti investimenti e acquisizioni di aziende del settore.
    \end{itemize}

    Inoltre entrambi sono state scelte per analizzare i trend del settore energetico nei periodo di tensioni geopolitiche, come la guerra in Ucraina e Iran e per il trend di calo durante la pandemia, essendo il settore energetico quello più colpito, visto le diminuzione di domanda del petrolio.

    \item Settore Automobilistico: 
    \begin{itemize}
        \item Tesla (TSLA): azienda specializzata nella produzione di veicoli elettrici (EV). È stata scelta poichè nonostante l'intensificarsi della concorrenza da parte dei costruttori cinesi, Tesla sta mantenendo una posizione di rilievo agli investimenti per la creazione di robotaxi.
        \item General Motors (GM): azienda leader in America basata su veicoli a motore. È stata scelta per il suo obbiettivo annunciato di produrre solo veicoli elettrici entro il 2035.
    \end{itemize}
\end{enumerate}

\section{Funzioni utilizzate}

\section{Breve Presentazione dei dati}



\chapter{Analisi Statistica}

\section{Rendimenti}

\section{Statistiche Descrittive}



\chapter{Previsione}

\section{Modello ARIMA}

\section{Strategie di Trading e Backtesting}

\section{Analisi Statica}

\section{Strategie Dinamiche}

\chapter{CAPM}

\chapter{Costruzione Portafoglio}

\end{document}
